\section{The physical chain, Fourier norms and dimension estimates}\label{sec:prelim}
This section fixes the common normalization. Every later passage from a sweep estimate to mixing uses physical time measured here. We use $n$ for a generic dyadic block size, including the full-deck size $N$ introduced above; the operators $T_n$ and $K_n$ are the corresponding sweep and reflected sweep on that block.

\subsection{Binary positions and sweep operators}
Number positions by $x=(x_1,\ldots,x_d)\in\mathbb F_2^d$, most significant bit first. One shuffle sends
\[
 x\longmapsto(x_2,\ldots,x_d,x_1+\xi_{x_2,\ldots,x_d}),
\]
where the $2^{d-1}$ bits $\xi$ are independent and fair. Write $R$ for the cyclic rotation and $h_t$ for the pair switches at physical time $t$, so the time-$t$ permutation is $Rh_t\cdots Rh_1$. Factoring out $R^t$ conjugates the switches into the successive coordinate directions. This deterministic left multiplication preserves distance to uniform. At time $d$, $R^d=I$, and one sweep has exactly the law of $d$ physical shuffles.

Permutations send each input position to its output. A product $gh$ applies $h$ first. For measures, $\mu*\nu$ denotes the law of $gh$ for independent $g\sim\mu,h\sim\nu$. In a unitary representation $\rho_\lambda$, define
\[
 \widehat\mu(\lambda)=\sum_g\mu(g)\rho_\lambda(g),
 \qquad\widehat{\mu*\nu}=\widehat\mu\,\widehat\nu.
\]
A fair switch layer averages the subgroup generated by its disjoint transpositions, and hence is an orthogonal projection $\Pi_i$. With chronological coordinate order $1,\ldots,d$, put
\[
 T_n=\Pi_d\cdots\Pi_1,\qquad K_n=T_n^*T_n,
 \qquad n=2^d.
\]
At the one-position recursive base, the empty product gives $T_1=K_1=I$.
\begin{lemma}[Annihilated shapes]\label{lem:annihilation}
The sign representation is annihilated by every nonempty fair layer. For $n\ge2$, every irreducible of $S_n$ indexed by a shape with more than $n/2$ rows is annihilated by a sweep.
\end{lemma}\begin{proof}
The sign average contains a fair transposition. For the second assertion, by Young's rule, the permutation module induced from the trivial representation of $(S_2)^{n/2}$ has only shapes dominating $(2^{n/2})$, hence at most $n/2$ rows.
\end{proof}

\begin{lemma}[Finite-size norm gap]\label{lem:finitegap}
For every fixed dyadic $n\ge2$, $\|T_n(\lambda)\|_{\op}<1$ on every nontrivial irreducible.
\end{lemma}
\begin{proof}
Equality on a unit vector would force equality at every orthogonal projection in the product. The vector would be fixed by all coordinate-pair transpositions. The edges of the cube form a connected graph, and its edge transpositions generate $S_n$, contradicting nontrivial irreducibility.
\end{proof}


\begin{lemma}[Support obstruction]\label{lem:support}
For every integer $t\ge0$,
\[
 \|q_d^{*t}-U_{S_n}\|_{\TV}\ge1-2^{tn/2}/n!.
\]
Consequently $t_{\mathrm{mix}}(d)\ge\lceil(2/n)\log_2(3n!/4)\rceil=2d-O(1)$.
\end{lemma}
\begin{proof}
At most $2^{tn/2}$ coin strings are available, so the law is supported on at most that many permutations. Comparing this support set with its uniform mass proves the first assertion. Distance at most $1/4$ requires support mass at least $3/4$, giving the exact integer lower bound. Stirling's formula gives the final expression.
\end{proof}

\subsection{Plancherel and powers of a nonnormal sweep}
All matrix traces and Schatten norms are unnormalized. Thus $\|A\|_{S^p}^p=\Tr|A|^p$, with $S^2=\HS$ and $S^\infty=\op$. If $D_\lambda$ is the irreducible dimension, finite-group orthogonality gives
\[
 |G|\sum_g|\mu(g)|^2=\sum_\lambda D_\lambda\|\widehat\mu(\lambda)\|_{\HS}^2.
\]
It applies to signed measures and subprobabilities as well as probabilities.
Together with Cauchy--Schwarz, it gives the finite-group Fourier upper
bound used by Diaconis and Shahshahani~\cite[Section 2, Lemma 2;
Section 3, Lemma 14]{DiaconisShahshahani1981}. For a measure $v$ of mass $m$,
\begin{equation}\label{eq:plancherel}
 \|v-mU_G\|_1^2\le
 \sum_{\lambda\ne\mathbf1}D_\lambda\|\widehat v(\lambda)\|_{\HS}^2.
\end{equation}
For a probability the left side is $4\|v-U_G\|_{\TV}^2$. If $0\le v\le\mu$ and $\mu$ is a probability, then
\begin{equation}\label{eq:lostmass}
 \|\mu-U_G\|_{\TV}\le1-m+\tfrac12\|v-mU_G\|_1.
\end{equation}
Both follow from Cauchy--Schwarz and the triangle inequality, respectively.

\begin{lemma}[Safe use of sweep powers]\label{lem:schatten}
For integers $r\ge s\ge1$ and any matrix $T$, setting $Q=T^*T$,
\[
 \|T^r\|_{\HS}^2\le\|Q\|_{\op}^{r-s}\Tr Q^s.
\]
\end{lemma}
\begin{proof}
Schatten H\"older gives $\|T^s\|_{\HS}\le\|T\|_{S^{2s}}^s$, and multiplying the remaining $r-s$ factors costs at most $\|T\|_{\op}^{r-s}$. Squaring proves the claim. No normality of $T$ is used.
\end{proof}

\subsection{Injection multiplicities and inverse-degree sums}
We use the standard indexing of complex irreducibles of $S_n$ by partitions $\lambda\vdash n$, with $D_\lambda=f^\lambda$ equal to the number of standard tableaux~\cite{Sagan2001,Stanley1999}. Write $k=n-\lambda_1$ and $\bar\lambda=(\lambda_2,\lambda_3,\ldots)$. Young branching and Frobenius reciprocity show that the representation on ordered distinct $r$-tuples contains $\lambda$ with multiplicity
\[
 f^{\lambda/(n-r)}.
\]
At $r=k$ this is $f^{\bar\lambda}$, and $\lambda$ does not occur on fewer slots. If $k\le r\le n-k$, the remaining first-row boxes are separated from the tail, giving
\begin{equation}\label{eq:tuplemultiplicity}
 m_\lambda(r)=\binom rk f^{\bar\lambda},\qquad
 D_\lambda\le\binom nk f^{\bar\lambda}.
\end{equation}
The inequality follows by choosing the entries below the first row and forgetting cross-row tableau constraints.
The same hook formula gives the useful quantitative refinement
\begin{equation}\label{eq:near-row-hooks}
 D_\lambda\ge\binom nk f^{\bar\lambda}
       \exp\left(-\frac{k}{n-2k+1}\right)
 \qquad(0\le k\le n/2).
\end{equation}
Indeed the hook inflation along the top row, relative to $(n-k)!$, is
\[
 \prod_{j\le\lambda_2}\left(1+
   \frac{\lambda'_j-1}{n-k-j+1}\right).
\]
The denominators are at least $n-2k+1$ and
$\sum_j(\lambda'_j-1)=k$. Taking logarithms bounds this product by
$\exp(k/(n-2k+1))$. The remaining hooks are exactly those of
$\bar\lambda$. In particular for $k\le.14n$ the loss is $\exp(O(k/n))$,
which also supplies the weaker $\exp(O(k\log n/n+k/n))$ loss when that
form is convenient.


The inverse-degree conclusions below are special cases of
Liebeck--Shalev~\cite[Proposition 2.5 and Theorem 2.6]{liebeck-shalev2004}.
We include an elementary proof of the estimates needed here, together
with a lower bound for individual dimensions.
\begin{lemma}[Degree sums]\label{lem:degrees}
Uniformly in $n$, $\sum_{\lambda\vdash n}D_\lambda^{-1}$ is bounded, and
\[
 \sum_{\lambda\ne(n),(1^n)}D_\lambda^{-1}\longrightarrow0.
\]
If $\ell=n-\max(\lambda_1,\lambda'_1)$, then
\[
 \log D_\lambda\ge(\log2)\sqrt\ell/2,
 \qquad \sup_n\sum_{\lambda\vdash n}D_\lambda^{-32}<\infty.
\]
\end{lemma}
\begin{proof}
Transpose if necessary so that the first row has length $r=\max(\lambda_1,\lambda'_1)$. If $r\le n/8$, the hook formula gives $D_\lambda\ge n!/(2r)^n$, exponential in $n$. If $n/8<r\le3n/4$, retain the first row and $\lfloor r/4\rfloor$ further boxes. Their tableaux extend to the whole diagram. A first row of length $r$ and a tail of size $j\le r$ have at least
\[
 \binom{r+j}{j}\frac{r-j+1}{r+1}
\]
tableaux: ballot interleavings of the first row with any fixed standard order of the tail respect all column inequalities. This is exponential in $n$ in the present range. There are $\exp(O(\sqrt n))$ partitions, so these ranges contribute $o(1)$ to the inverse-degree sum. If $r>3n/4$, write $j=n-r$. The same ballot bound is at least a fixed multiple of $\binom nj$. There are at most $2p(j)$ possible shapes up to transpose, and $\binom nj\ge4^j$ for $j<n/4$. Since $p(j)=\exp(O(\sqrt j))$ and each fixed positive $j$ gives a divergent binomial coefficient, dominated convergence proves the first assertion.

For the square-root estimate, put $b=\lambda_2$ and $h=\lambda'_1$, so $b(h-1)\ge\ell$. If $h-1\ge\sqrt\ell$, a hook with row and column length $h$ has at least $2^{h-1}$ tableaux. Otherwise $b\ge\sqrt\ell$ and the two-row rectangle of width $b$ has the Catalan number of tableaux, at least $2^{b-1}$. Subdiagram tableaux extend; handling $\ell=0$ separately proves the displayed bound. Finally $p(j)\le e^{3\sqrt j}$ follows by evaluating the partition generating product at $e^{-1/\sqrt j}$. Summing $2e^{3\sqrt\ell}e^{-16(\log2)\sqrt\ell}$ proves the inverse-32 bound.
\end{proof}
